\begin{abstract}
    \noindent
    O crescimento do volume de dados em \textit{data lakes} tem impulsionado a necessidade de arquiteturas adaptativas que equilibrem custo e \textit{performance} na execução de consultas SQL. Este trabalho apresenta o \textit{DyraSQL}, um \textit{framework} para roteamento dinâmico de consultas SQL baseado em metadados de tabelas \textit{Apache Iceberg}, obtidos por meio do comando \textit{EXPLAIN (TYPE IO)} do \textit{Trino}. O sistema direciona automaticamente consultas para \textit{clusters} de processamento com capacidades distintas, classificados como \textit{small}, \textit{medium} e \textit{large}, a partir de um algoritmo de pontuação que combina volume de dados, complexidade da consulta e histórico de execuções. A metodologia envolve extração de metadados sobre volume efetivo e custos de I/O, análise sintática da consulta SQL e reaproveitamento de decisões em \textit{cache} temporal no PostgreSQL. O protótipo foi avaliado em ambiente local orquestrado por Docker Compose, com tabelas \textit{Apache Iceberg} geradas sobre armazenamento de objetos MinIO e catálogo REST, demonstrando que o sistema calcula \textit{scores} consistentes e classifica consultas nas faixas de capacidade definidas. Os resultados indicam que o roteamento baseado em metadados contribui para o uso mais eficiente de recursos de processamento, oferecendo um modelo adaptativo aplicável a \textit{clusters} \textit{Trino} locais e, em trabalhos futuros, a ambientes de nuvem.
    \\\textbf{\keyword{Palavras-chave: }} Consultas SQL. Trino. Apache Iceberg. Data lake. Trino Gateway.
    \end{abstract}
